% ------------------------------------------------------------
% VENTAJAS COMPARATIVAS: Binance/BTC vs NASA HTTP
% ------------------------------------------------------------

\begin{itemize}[leftmargin=2em]
    \item \textbf{Relevancia temporal.}
          El dataset NASA HTTP corresponde a julio--agosto de 1995, con
          arquitecturas web, patrones de usuarios y tecnologías de red
          completamente obsoletas. Los datos de Bitcoin son contemporáneos
          y reflejan dinámicas de mercado actuales, aumentando la validez
          externa del análisis.

    \item \textbf{Regularidad muestral.}
          La API de Binance entrega observaciones equiespaciadas en el
          tiempo (velas 1m/5m), mientras que los tiempos entre llegadas
          HTTP son irregulares por naturaleza. La regularidad facilita la
          aplicación directa de estimadores fractales (VT, LD, wavelet
          leaders) sin necesidad de resampling o interpolación que
          introduzca artefactos.

    \item \textbf{Reproducibilidad y extensibilidad.}
          La API pública de Binance permite a cualquier lector del trabajo
          replicar exactamente la serie descargada, o extenderla a otros
          pares (ETH/USDT, BNB/USDT), intervalos (5m, 1h, 1d) o ventanas
          temporales. El dataset NASA, por el contrario, es un archivo
          estático de dominio público cuyo contenido no puede regenerarse.

    \item \textbf{Fundamentación teórica robusta.}
          Existe literatura indexada abundante sobre LRD en criptomonedas
          (Bariviera et al., 2017; Dionisio et al., 2022), lo que permite
          contrastar los resultados obtenidos contra valores de referencia
          establecidos. La literatura sobre fractalidad en tráfico web de
          1995 es mucho más limitada y de contexto histórico específico.

    \item \textbf{Narrativa transversal fortalecida.}
          Bitcoin ($H\approx0{,}6$--$0{,}7$) ocupa una posición
          intermedia en la taxonomía entre NASA HTTP ($H\approx0{,}54$)
          y corazón sano ($H\approx0{,}96$), transformando el eje
          horizontal del mapa conceptual en un gradiente continuo de
          memoria persistente en lugar de un salto abrupto entre
          tecnología y biología.

    \item \textbf{Mercado 24/7 sin gaps.}
          La ausencia de cierres diarios elimina discontinuidades
          estructurales que en mercados tradicionales requieren corrección
          (ajustes por overnight risk, efectos weekend), simplificando el
          preprocesamiento y la interpretación de los estimadores.
\end{itemize}
